% GENERATED FILE. Edit canonical lessons, then run npm run generate:books.
% canonical-source-hash: fnv1a64:da8a7f6cc1c74f94
% canonical-lessons: SA-C347-udaharanam, SA-C347-drstantah, SA-C347-ittham, SA-C347-bodhayati, SA-C347-abhiprayah, SA-R347-review-340-347

\chapter{For Example}
\label{ch:sa-for-example}
\hlchaptermodality{347}

\begin{chapteropening}
\textbf{By the end of this chapter:} \emph{I can back up a point with an example.}

Complete the last lesson of chapter 347: I can back up a point with an example.
\end{chapteropening}

\section[\texorpdfstring{\sk{उदाहरणम्} (udāharaṇam)}{उदाहरणम् (udāharaṇam)}]{\sk{उदाहरणम्} (udāharaṇam) — an example}

\label{lesson:SA-C347-udaharanam}



\begin{quote}
Before the new one: say the Sanskrit for a doubt, then the Sanskrit for in fact.
\end{quote}

\subsection*{You'll want to know: \sk{उदाहरणम्}}
\textbf{\sk{उदाहरणम्}} — \emph{udāharaṇam} — \textquotedblleft{}an example\textquotedblright{}.

\begin{quote}
\sk{कृपया} \sk{एकम्} \sk{उदाहरणं} \sk{ददातु।}
\end{quote}

\begin{quote}
\emph{kṛpayā ekam udāharaṇaṃ dadātu.}
\end{quote}

\begin{quote}
Please give one example.
\end{quote}

A neuter noun.

\begin{cousinweb}
\textbf{\sk{उदाहरणम्}} (\emph{udāharaṇam}) is built on \emph{ud}, \emph{ā} and the root √hṛ, \textquotedblleft{}to take, to carry\textquotedblright{}; MW gives \emph{ud-āharaṇa} \textquotedblleft{}declaration\textquotedblright{} and \textquotedblleft{}an example, an illustration\textquotedblright{} (MW).
\end{cousinweb}

\begin{grammarlens}[title={What you've built}]
One more word for giving reasons and opinions.
\end{grammarlens}

\subsection*{Your turn}
\begin{itemize}
\raggedright
  \item \emph{Say it:} \emph{udāharaṇam}
  \item \emph{Say it:} \emph{udāharaṇam}, once more
  \item \emph{Say it:} \emph{kṛpayā ekam udāharaṇaṃ dadātu.} — Please give one example.
  \item \emph{Recall:} say the Sanskrit for a doubt, then the Sanskrit for in fact, then say \emph{udāharaṇam} again
\end{itemize}

\begin{tcolorbox}[breakable,colback=teal!4,colframe=teal!35!black,title={Before you move on}]
What does \sk{उदाहरणम्} mean? (\textquotedblleft{}An example\textquotedblright{}.) Say it once more.
\end{tcolorbox}
\section[\texorpdfstring{\sk{दृष्टान्तः} (dṛṣṭāntaḥ)}{दृष्टान्तः (dṛṣṭāntaḥ)}]{\sk{दृष्टान्तः} (dṛṣṭāntaḥ) — an illustration (in an argument)}

\label{lesson:SA-C347-drstantah}



\begin{quote}
Before the new one: say the Sanskrit for in fact, then the Sanskrit for an example.
\end{quote}

\subsection*{You'll want to know: \sk{दृष्टान्तः}}
\textbf{\sk{दृष्टान्तः}} — \emph{dṛṣṭāntaḥ} — \textquotedblleft{}an illustration (in an argument)\textquotedblright{}.

\begin{quote}
\sk{गुरुः} \sk{दृष्टान्तेन} \sk{विचारं} \sk{स्पष्टं} \sk{करोति।}
\end{quote}

\begin{quote}
\emph{guruḥ dṛṣṭāntena vicāraṃ spaṣṭaṃ karoti.}
\end{quote}

\begin{quote}
The teacher makes the idea clear with an illustration.
\end{quote}

A masculine noun.

\begin{cousinweb}
\textbf{\sk{दृष्टान्तः}} (\emph{dṛṣṭāntaḥ}) is built from \emph{dṛṣṭa}, \textquotedblleft{}seen\textquotedblright{}, and \emph{anta}, \textquotedblleft{}an end\textquotedblright{}, which MW explains as the end or aim of what is seen; MW gives the senses \textquotedblleft{}example, standard\textquotedblright{} and \textquotedblleft{}exemplification\textquotedblright{} (MW).
\end{cousinweb}

\begin{grammarlens}[title={What you've built}]
One more word for giving reasons and opinions.
\end{grammarlens}

\subsection*{Your turn}
\begin{itemize}
\raggedright
  \item \emph{Say it:} \emph{dṛṣṭāntaḥ}
  \item \emph{Say it:} \emph{dṛṣṭāntaḥ}, once more
  \item \emph{Say it:} \emph{guruḥ dṛṣṭāntena vicāraṃ spaṣṭaṃ karoti.} — The teacher makes the idea clear with an illustration.
  \item \emph{Recall:} say the Sanskrit for in fact, then the Sanskrit for an example, then say \emph{dṛṣṭāntaḥ} again
\end{itemize}

\begin{tcolorbox}[breakable,colback=teal!4,colframe=teal!35!black,title={Before you move on}]
What does \sk{दृष्टान्तः} mean? (\textquotedblleft{}An illustration (in an argument)\textquotedblright{}.) Say it once more.
\end{tcolorbox}
\section[\texorpdfstring{\sk{इत्थम्} (ittham)}{इत्थम् (ittham)}]{\sk{इत्थम्} (ittham) — in this way; thus}

\label{lesson:SA-C347-ittham}



\begin{quote}
Before the new one: say the Sanskrit for an example, then the Sanskrit for an illustration (in an argument).
\end{quote}

\subsection*{You'll want to know: \sk{इत्थम्}}
\textbf{\sk{इत्थम्}} — \emph{ittham} — \textquotedblleft{}in this way; thus\textquotedblright{}.

\begin{quote}
\sk{भवान्} \sk{इत्थं} \sk{करोतु।}
\end{quote}

\begin{quote}
\emph{bhavān itthaṃ karotu.}
\end{quote}

\begin{quote}
Please do it this way.
\end{quote}

\begin{cousinweb}
\textbf{\sk{इत्थम्}} (\emph{ittham}) is built on the pronoun stem \emph{id}, \textquotedblleft{}this\textquotedblright{}; MW gives \emph{ittham} \textquotedblleft{}thus, in this manner\textquotedblright{} (MW).
\end{cousinweb}

\begin{grammarlens}[title={What you've built}]
One more word for giving reasons and opinions.
\end{grammarlens}

\subsection*{Your turn}
\begin{itemize}
\raggedright
  \item \emph{Say it:} \emph{ittham}
  \item \emph{Say it:} \emph{ittham}, once more
  \item \emph{Say it:} \emph{bhavān itthaṃ karotu.} — Please do it this way.
  \item \emph{Recall:} say the Sanskrit for an example, then the Sanskrit for an illustration (in an argument), then say \emph{ittham} again
\end{itemize}

\begin{tcolorbox}[breakable,colback=teal!4,colframe=teal!35!black,title={Before you move on}]
What does \sk{इत्थम्} mean? (\textquotedblleft{}In this way; thus\textquotedblright{}.) Say it once more.
\end{tcolorbox}
\section[\texorpdfstring{\sk{बोधयति} (bodhayati)}{बोधयति (bodhayati)}]{\sk{बोधयति} (bodhayati) — makes (someone) understand}

\label{lesson:SA-C347-bodhayati}



\begin{quote}
Before the new one: say the Sanskrit for an illustration (in an argument), then the Sanskrit for in this way.
\end{quote}

\subsection*{You'll want to know: \sk{बोधयति}}
\textbf{\sk{बोधयति}} — \emph{bodhayati} — \textquotedblleft{}makes (someone) understand\textquotedblright{}.

\begin{quote}
\sk{अध्यापकः} \sk{उदाहरणेन} \sk{छात्रान्} \sk{बोधयति।}
\end{quote}

\begin{quote}
\emph{adhyāpakaḥ udāharaṇena chātrān bodhayati.}
\end{quote}

\begin{quote}
The teacher makes the students understand with an example.
\end{quote}

Past (he …): \sk{सः} \sk{छात्रान्} \sk{बोधितवान्}, \emph{saḥ chātrān bodhitavān}.

\begin{cousinweb}
\textbf{\sk{बोधयति}} (\emph{bodhayati}) comes from the root √budh, \textquotedblleft{}to wake, to understand\textquotedblright{}, in its causative form; MW gives \emph{bodhita} \textquotedblleft{}made known, apprised, explained\textquotedblright{} (MW).
\end{cousinweb}

\begin{grammarlens}[title={What you've built}]
One more word for giving reasons and opinions.
\end{grammarlens}

\subsection*{Your turn}
\begin{itemize}
\raggedright
  \item \emph{Say it:} \emph{bodhayati}
  \item \emph{Say it:} \emph{bodhayati}, once more
  \item \emph{Say it:} \emph{adhyāpakaḥ udāharaṇena chātrān bodhayati.} — The teacher makes the students understand with an example.
  \item \emph{Recall:} say the Sanskrit for an illustration (in an argument), then the Sanskrit for in this way, then say \emph{bodhayati} again
\end{itemize}

\begin{tcolorbox}[breakable,colback=teal!4,colframe=teal!35!black,title={Before you move on}]
What does \sk{बोधयति} mean? (\textquotedblleft{}Makes (someone) understand\textquotedblright{}.) Say it once more.
\end{tcolorbox}
\section[\texorpdfstring{\sk{अभिप्रायः} (abhiprāyaḥ)}{अभिप्रायः (abhiprāyaḥ)}]{\sk{अभिप्रायः} (abhiprāyaḥ) — what someone means; an intention}

\label{lesson:SA-C347-abhiprayah}



\begin{quote}
Before the new one: say the Sanskrit for in this way, then the Sanskrit for makes (someone) understand.
\end{quote}

\subsection*{You'll want to know: \sk{अभिप्रायः}}
\textbf{\sk{अभिप्रायः}} — \emph{abhiprāyaḥ} — \textquotedblleft{}what someone means; an intention\textquotedblright{}.

\begin{quote}
\sk{तव} \sk{अभिप्रायः} \sk{कः}?
\end{quote}

\begin{quote}
\emph{tava abhiprāyaḥ kaḥ?}
\end{quote}

\begin{quote}
What do you mean?
\end{quote}

A masculine noun.

\begin{cousinweb}
\textbf{\sk{अभिप्रायः}} (\emph{abhiprāyaḥ}) is built on \emph{abhi}, \emph{pra} and the root √i, \textquotedblleft{}to go\textquotedblright{}; MW gives \emph{abhi-prāya} \textquotedblleft{}aim, purpose, intention\textquotedblright{} and \textquotedblleft{}meaning, sense\textquotedblright{} (MW).
\end{cousinweb}

\begin{grammarlens}[title={What you've built}]
One more word for giving reasons and opinions.
\end{grammarlens}

\subsection*{Your turn}
\begin{itemize}
\raggedright
  \item \emph{Say it:} \emph{abhiprāyaḥ}
  \item \emph{Say it:} \emph{abhiprāyaḥ}, once more
  \item \emph{Say it:} \emph{tava abhiprāyaḥ kaḥ?} — What do you mean?
  \item \emph{Recall:} say the Sanskrit for in this way, then the Sanskrit for makes (someone) understand, then say \emph{abhiprāyaḥ} again
\end{itemize}

\begin{tcolorbox}[breakable,colback=teal!4,colframe=teal!35!black,title={Before you move on}]
What does \sk{अभिप्रायः} mean? (\textquotedblleft{}What someone means; an intention\textquotedblright{}.) Say it once more.
\end{tcolorbox}
\section[(dialogue)]{Looking back: from Why It Happened to For Example}

\label{lesson:SA-R347-review-340-347}



\begin{quote}
Say the Sanskrit for an occasion and an influence.
\end{quote}

\subsection*{Your turn}
Say these aloud:

\begin{itemize}
\raggedright
  \item an occasion, an influence, a lack, comes about, gives rise to
  \item a purpose, an aim, a use, urges on, has something done
  \item a basis, a line of reasoning, that is to say, as a result, argues for
  \item a debate, reasoning, a speaker, a listener, sets out (a case)
  \item the opposing side, opposition, supports (a view), consents, is it not so
  \item conviction, testimony, a fact, proves, certainty
  \item an objection, refutes, rejects (a claim), a doubt, in fact
  \item an example, an illustration (in an argument), in this way, makes (someone) understand, what someone means
\end{itemize}

\begin{tcolorbox}[breakable,colback=teal!4,colframe=teal!35!black,title={Before you move on}]
An occasion? (\textbf{\sk{निमित्तम्}}, \emph{nimittam}.) An influence? (\textbf{\sk{प्रभावः}}, \emph{prabhāvaḥ}.) A lack? (\textbf{\sk{अभावः}}, \emph{abhāvaḥ}.) Comes about? (\textbf{\sk{उत्पद्यते}}, \emph{utpadyate}.) Gives rise to? (\textbf{\sk{जनयति}}, \emph{janayati}.) A purpose? (\textbf{\sk{प्रयोजनम्}}, \emph{prayojanam}.) An aim? (\textbf{\sk{उद्देश्यम्}}, \emph{uddeśyam}.) A use? (\textbf{\sk{उपयोगः}}, \emph{upayogaḥ}.) Urges on? (\textbf{\sk{प्रेरयति}}, \emph{prerayati}.) Has something done? (\textbf{\sk{कारयति}}, \emph{kārayati}.) A basis? (\textbf{\sk{आधारः}}, \emph{ādhāraḥ}.) A line of reasoning? (\textbf{\sk{युक्तिः}}, \emph{yuktiḥ}.) That is to say? (\textbf{\sk{अर्थात्}}, \emph{arthāt}.) As a result? (\textbf{\sk{फलतः}}, \emph{phalataḥ}.) Argues for? (\textbf{\sk{उपपादयति}}, \emph{upapādayati}.) A debate? (\textbf{\sk{वादः}}, \emph{vādaḥ}.) Reasoning? (\textbf{\sk{तर्कः}}, \emph{tarkaḥ}.) A speaker? (\textbf{\sk{वक्ता}}, \emph{vaktā}.) A listener? (\textbf{\sk{श्रोता}}, \emph{śrotā}.) Sets out (a case)? (\textbf{\sk{प्रतिपादयति}}, \emph{pratipādayati}.) The opposing side? (\textbf{\sk{प्रतिपक्षः}}, \emph{pratipakṣaḥ}.) Opposition? (\textbf{\sk{विरोधः}}, \emph{virodhaḥ}.) Supports (a view)? (\textbf{\sk{समर्थयति}}, \emph{samarthayati}.) Consents? (\textbf{\sk{अनुमन्यते}}, \emph{anumanyate}.) Is it not so? (\textbf{\sk{ननु}}, \emph{nanu}.) Conviction? (\textbf{\sk{प्रत्ययः}}, \emph{pratyayaḥ}.) Testimony? (\textbf{\sk{साक्ष्यम्}}, \emph{sākṣyam}.) A fact? (\textbf{\sk{तथ्यम्}}, \emph{tathyam}.) Proves? (\textbf{\sk{साधयति}}, \emph{sādhayati}.) Certainty? (\textbf{\sk{निश्चयः}}, \emph{niścayaḥ}.) An objection? (\textbf{\sk{आक्षेपः}}, \emph{ākṣepaḥ}.) Refutes? (\textbf{\sk{खण्डयति}}, \emph{khaṇḍayati}.) Rejects (a claim)? (\textbf{\sk{निरस्यति}}, \emph{nirasyati}.) A doubt? (\textbf{\sk{संशयः}}, \emph{saṃśayaḥ}.) In fact? (\textbf{\sk{वस्तुतः}}, \emph{vastutaḥ}.) An example? (\textbf{\sk{उदाहरणम्}}, \emph{udāharaṇam}.) An illustration (in an argument)? (\textbf{\sk{दृष्टान्तः}}, \emph{dṛṣṭāntaḥ}.) In this way? (\textbf{\sk{इत्थम्}}, \emph{ittham}.) Makes (someone) understand? (\textbf{\sk{बोधयति}}, \emph{bodhayati}.) What someone means? (\textbf{\sk{अभिप्रायः}}, \emph{abhiprāyaḥ}.)
\end{tcolorbox}
